\section{Connected components}
\label{section-connected-components}

\begin{definition}
\label{definition-connected-components}
Let $X$ be a topological space.
\begin{enumerate}
\item We say $X$ is {\it connected} if $X$ is not empty and whenever
$X = T_1 \amalg T_2$ with $T_i \subset X$ open and closed, then either
$T_1 = \emptyset$ or $T_2 = \emptyset$.
\item We say $T \subset X$ is a {\it connected component} of $X$ if
$T$ is a maximal connected subset of $X$.
\end{enumerate}
\end{definition}

\noindent
The empty space is not connected.

\begin{lemma}
\label{lemma-image-connected-space}
Let $f : X \to Y$ be a continuous map of topological spaces.
If $E \subset X$ is a connected subset, then $f(E) \subset Y$
is connected as well.
\end{lemma}

\begin{proof}
Let $A \subset f(E)$ an open and closed subset of $f(E)$. Because $f$ is
continuous, $f^{-1}(A)$ is an open and closed subset of $E$. As $E$ is
connected, $f^{-1}(A) = \emptyset$ or $f^{-1}(A) = E$. However, $A\subset f(E)$
implies that $A = f(f^{-1}(A))$. Indeed, if $x\in f(f^{-1}(A))$ then there is
$y\in f^{-1}(A)$ such that $f(y) = x$ and because $y\in f^{-1}(A)$, we have
$f(y) \in A$ i.e. $x\in A$. Reciprocally, if $x\in A$, $A\subset f(E)$ implies
that there is $y\in E$ such that $f(y) = x$. Therefore $y\in f^{-1}(A)$, and
then $x\in f(f^{-1}(A))$. Thus $A = \emptyset$ or $A = f(E)$ proving that
$f(E)$ is connected.
\end{proof}

\begin{lemma}
\label{lemma-connected-components}
Let $X$ be a topological space.
\begin{enumerate}
\item If $T \subset X$ is connected, then so is its closure.
\item Any connected component of $X$ is closed (but not necessarily open).
\item Every connected subset of $X$ is contained in a unique connected
component of $X$.
\item Every point of $X$ is contained in a unique connected component, in other
words, $X$ is the disjoint union of its connected components.
\end{enumerate}
\end{lemma}

\begin{proof}
Let $\overline{T}$ be the closure of the connected subset $T$.
Suppose $\overline{T} = T_1 \amalg T_2$ with $T_i \subset \overline{T}$
open and closed. Then $T = (T\cap T_1) \amalg (T \cap T_2)$. Hence
$T$ equals one of the two, say $T = T_1 \cap T$. Thus
$\overline{T} \subset T_1$. This implies (1) and (2).

\medskip\noindent
Let $A$ be a nonempty set of connected subsets of $X$ such that
$\Omega = \bigcap_{T \in A} T$ is nonempty. We claim $E = \bigcup_{T \in A} T$
is connected. Namely, $E$ is nonempty as it contains $\Omega$.
Say $E = E_1 \amalg E_2$ with $E_i$ closed in $E$. We may assume
$E_1$ meets $\Omega$ (after renumbering). Then each $T \in A$ meets $E_1$
and hence must be contained in $E_1$ as $T$ is connected. Hence $E \subset E_1$
which proves the claim.

\medskip\noindent
Let $W \subset X$ be a nonempty connected subset.
If we apply the result of the previous paragraph to the set of all
connected subsets of $X$ containing $W$, then we see that $E$ is
a connected component of $X$. This implies existence and uniqueness in (3).

\medskip\noindent
Let $x \in X$. Taking $W = \{x\}$ in the previous paragraph
we see that $x$ is contained in a unique connected component of $X$.
Any two distinct connected components must be disjoint (by the
result of the second paragraph).

\medskip\noindent
To get an example where connected components are not open, just take
an infinite product $\prod_{n \in \mathbf{N}} \{0, 1\}$
with the product topology. Its connected components are singletons,
which are not open.
\end{proof}

\begin{remark}
\label{remark-quasi-components}
\begin{reference}
\cite[Example 6.1.24]{Engelking}
\end{reference}
Let $X$ be a topological space and $x \in X$. Let $Z \subset X$ be the
connected component of $X$ passing through $x$. Consider the intersection
$E$ of all open and closed subsets of $X$ containing $x$. It is clear that
$Z \subset E$. In general $Z \not = E$. For example, let
$X = \{x, y, z_1, z_2, \ldots\}$ with the topology with the following
basis of opens, $\{z_n\}$, $\{x, z_n, z_{n + 1}, \ldots\}$, and
$\{y, z_n, z_{n + 1}, \ldots\}$ for all $n$. Then $Z = \{x\}$ and
$E = \{x, y\}$. We omit the details.
\end{remark}

\begin{lemma}
\label{lemma-connected-fibres-quotient-topology-connected-components}
Let $f : X \to Y$ be a continuous map of topological spaces.
Assume that
\begin{enumerate}
\item all fibres of $f$ are connected, and
\item a set $T \subset Y$ is closed if and only if $f^{-1}(T)$ is closed.
\end{enumerate}
Then $f$ induces a bijection between the sets of connected
components of $X$ and $Y$.
\end{lemma}

\begin{proof}
Let $T \subset Y$ be a connected component.
Note that $T$ is closed, see Lemma \ref{lemma-connected-components}.
The lemma follows if we show that $f^{-1}(T)$ is connected
because any connected subset of $X$ maps into a connected component
of $Y$ by Lemma \ref{lemma-image-connected-space}.
Suppose that $f^{-1}(T) = Z_1 \amalg Z_2$
with $Z_1$, $Z_2$ closed. For any $t \in T$ we see that
$f^{-1}(\{t\}) = Z_1 \cap f^{-1}(\{t\}) \amalg Z_2 \cap f^{-1}(\{t\})$.
By (1) we see $f^{-1}(\{t\})$ is connected we conclude that
either $f^{-1}(\{t\}) \subset Z_1$ or $f^{-1}(\{t\}) \subset Z_2$.
In other words $T = T_1 \amalg T_2$ with $f^{-1}(T_i) = Z_i$.
By (2) we conclude that $T_i$ is closed in $Y$.
Hence either $T_1 = \emptyset$ or $T_2 = \emptyset$ as desired.
\end{proof}

\begin{lemma}
\label{lemma-connected-fibres-connected-components}
Let $f : X \to Y$ be a continuous map of topological spaces.
Assume that
(a) $f$ is open,
(b) all fibres of $f$ are connected.
Then $f$ induces a bijection between the sets of connected
components of $X$ and $Y$.
\end{lemma}

\begin{proof}
This is a special case of
Lemma \ref{lemma-connected-fibres-quotient-topology-connected-components}.
\end{proof}

\begin{lemma}
\label{lemma-finite-fibre-connected-components}
Let $f : X \to Y$ be a continuous map of nonempty topological spaces. Assume
that
(a) $Y$ is connected,
(b) $f$ is open and closed, and
(c) there is a point $y\in Y$ such that the fiber $f^{-1}(y)$ is a finite set.
Then $X$ has at most $|f^{-1}(y)|$ connected components. Hence any connected 
component $T$ of $X$ is open and closed, and $f(T)$ is a nonempty open and 
closed subset of $Y$, which is therefore equal to $Y$.
\end{lemma}

\begin{proof}
If the topological space $X$ has at least $N$ connected components for some
$N \in \mathbf{N}$, we find by induction a decomposition
$X = X_1 \amalg \ldots \amalg X_N$ of $X$ as a disjoint union of $N$ nonempty
open and closed subsets $X_1, \ldots , X_N$ of $X$. As $f$ is open and closed,
each $f(X_i)$ is a nonempty open and closed subset of $Y$ and is hence equal to
$Y$. In particular the intersection $X_i \cap f^{-1}(y)$ is nonempty for each
$1 \leq i \leq N$. Hence $f^{-1}(y)$ has at least $N$ elements.
\end{proof}

\begin{definition}
\label{definition-totally-disconnected}
A topological space is {\it totally disconnected} if the connected components
are all singletons.
\end{definition}

\noindent
A discrete space is totally disconnected.
A totally disconnected space need not be discrete, for example
$\mathbf{Q} \subset \mathbf{R}$ is totally disconnected but not discrete.

\begin{lemma}
\label{lemma-space-connected-components}
Let $X$ be a topological space. Let $\pi_0(X)$ be the set of connected
components of $X$. Let $X \to \pi_0(X)$ be the map which sends
$x \in X$ to the connected component of $X$ passing through $x$.
Endow $\pi_0(X)$ with the quotient topology. Then $\pi_0(X)$ is a
totally disconnected space and any continuous map $X \to Y$
from $X$ to a totally disconnected space $Y$ factors through $\pi_0(X)$.
\end{lemma}

\begin{proof}
By Lemma
\ref{lemma-connected-fibres-quotient-topology-connected-components}
the connected components of $\pi_0(X)$ are the singletons.
We omit the proof of the second statement.
\end{proof}

\begin{definition}
\label{definition-locally-connected}
A topological space $X$ is called {\it locally connected} if
every point $x \in X$ has a fundamental system of connected neighbourhoods.
\end{definition}

\begin{lemma}
\label{lemma-locally-connected}
Let $X$ be a topological space. If $X$ is locally connected, then
\begin{enumerate}
\item any open subset of $X$ is locally connected, and
\item the connected components of $X$ are open.
\end{enumerate}
So also the connected components of open subsets of $X$ are open.
In particular, every point has a fundamental system of open connected
neighbourhoods.
\end{lemma}

\begin{proof}
For all $x\in X$ let write $\mathcal{N}(x)$ the fundamental system of connected
neighbourhoods of $x$ and let $U\subset X$ be an open subset of $X$. Then for
all $x\in U$, $U$ is a neighbourhood of $x$, so the set
$\{V \in \mathcal{N}(x) | V\subset U\}$
is not empty and is a fundamental system of connected
neighbourhoods of $x$ in $U$. Thus $U$ is locally connected and it proves (1).

\medskip\noindent
Let $x \in \mathcal{C} \subset X$ where $\mathcal{C}$ is the connected component
of $x$. Because $X$ is locally connected, there exists $\mathcal{N}$ a
connected neighbourhood of $x$. Therefore by the definition of a connected
component, we have $\mathcal{N} \subset \mathcal{C}$ and then $\mathcal{C}$ is
a neighbourhood of $x$. It implies that $\mathcal{C}$ is a neighbourhood of
each of his point, in other words $\mathcal{C}$ is open and (2) is proven.
\end{proof}




